\section{Hole-targeted values, stock readers, and cost}
\label{sec:appendix:systems}

\paragraph{Hole-targeted set}
\path{pipeline/adaptive.py} builds, per host language, values aimed at the
hole rather than at a native closer. A grammar puts one to three hole-boundary
atoms before a breakout core and up to two after it: matchers, quotes, the
matchertext escapes, comments, a zero byte, and URL-encoded and fullwidth
matchers; for SQL also the overlong UTF-8 forms of \verb|)| and \verb|]|, and
for HTML also \verb|]>|, end tags, comment delimiters, and a character whose
last UTF-8 byte Shift\_JIS reads as a lead byte. Each grammar value enters as
built and made matchertext by added matchers, which \textsc{Verify} accepts. An
even sample of up to 500 effective payloads enters with eleven variants each:
made matchertext; led or trailed by each lone closer; and made matchertext, then
wrapped in \verb|()| or \verb|[]|, with each matcher backslash-escaped, or with a
trailing zero byte. Two cores enter nested at the depth limit and one past it.
Each value is canonicalized (A3), and values of the recorded corpus are left
out. For SQL, \dataSqliTargetedValues values break out of the controls in
\dataSqliTargetedEffective (value, fixture) pairs and of the holes in none of
\dataSqliTargetedHoleCases cases. For HTML, \dataHtmlTargetedValues values break
out of concatenation in \dataHtmlTargetedEffective (value, host) pairs and of no
matchertext arm or guarded writer.

\paragraph{Charset}
Two more HTML arms send the verified-hole document of each host as UTF-8 bytes.
Decoded by the WHATWG UTF-8 decoder, all \dataHtmlUtfEightRead such documents of
the hole-targeted set parse as the string does. Decoded as Shift\_JIS,
\dataHtmlSjisBreakout break out. In the value \={a}\verb|[x]><img src=x onerror=alert(1)>|,
the last UTF-8 byte of \={a} (U+0101) is a Shift\_JIS lead byte, which takes the
\verb|[| as its second byte, so the reader ends the hole at the \verb|]| that
\textsc{Verify} matched with it. The decoding of A3 thus starts at the bytes: the
writer declares the charset it verified under.

\paragraph{Stock readers}
\path{test/matchertext/csrc/stock_reader.c} is SQLite built from the same tree
without matchertext. It reads the statements that the text holes compose for the
corpus and the hole-targeted set. It refuses all \dataSqliDowngradeValueCases
value-hole statements, since it reads \verb|M'(| as a name before a string. Of
\dataSqliDowngradeNameCases name-hole statements, \dataSqliDowngradeNameBreakout
break out: the encoder writes the value \verb|]AND 0/*M'()| as
\verb|[\c[]AND 0/*M'()]|, and stock SQLite ends the name at the \verb|]| of
\verb|\c[]|.

\paragraph{Stock SQL}
The additive build passes every case of SQLite's full Tcl test suite
(\dataSqliteSuiteTests cases) that the stock build passes; the stock build
itself fails \dataSqliteSuiteStockErrors cases on our test host. Three forms that
the suite lacks read differently. A bracketed name with an unmatched matcher is
an error (\verb|[a(b]|), a bracketed name with a matchertext escape is decoded
(\verb|[a\c()b]| names \verb|a)b|), and an \verb|m| directly before a string that
opens with a matcher is a value hole (\verb|SELECT m'(x)'| returns \texttt{x},
not column \texttt{m} named \texttt{(x)}).

\paragraph{Stock HTML}
With both matchertext options off, the parse5 fork passes parse5's test suite
unchanged. The serializer guard, on by default, changes two html5lib serializer
tests, both trees whose stock output parses back as another tree: it refuses
\texttt{plaintext}, which no end tag closes, and it writes back the line feed
that a parser drops at the start of \texttt{pre}, \texttt{textarea}, and
\texttt{listing}. With the parser option on, the html5lib tokenizer and tree
construction tests pass unchanged, and Angular's \verb|[prop]="x"| attributes
parse as in stock parse5. One stock form reads differently: an attribute name
in brackets placed last with no value (\verb|<div [x]>|) is an element hole.

\paragraph{Cost}
\path{pipeline/perf.py} reports medians on one machine. The additive build
runs \texttt{speedtest1} (main test set, in memory) in \dataPerfSpeedtestRatio of
the stock time. A typed \texttt{?V} prepare costs \dataPerfTypedShortExtraNs\,ns
more than a prepare and bind of a \dataPerfTypedShortLength-byte value, and
\dataPerfTypedNsPerByte\,ns more per byte of a longer value; the in-band
\verb|M'(|$\dots$\verb|)'| hole costs \dataPerfMtNsPerByte\,ns per byte. The
parse5 fork parses its benchmark page in \dataPerfParseRatio of the time of its
upstream base, with the option off or on, and serializes it in
\dataPerfSerializeRatio with the guard on.

\section{HTML in browsers, nested layers, and mutation XSS}
\label{sec:appendix:html}

\paragraph{Browsers}
\path{pipeline/browsers.py} installs the Chromium \dataHtmlChromiumVersion, Firefox
\dataHtmlFirefoxVersion, and WebKit \dataHtmlWebkitVersion builds that Playwright
\dataPlaywrightVersion pins~\cite{playwright}; they are not branded Chrome,
Firefox, or Safari, and Playwright picks the WebKit build by operating system.
Each browser reads a document as a page load does: a document that contains
\texttt{noscript}, the one element that depends on scripting, loads in a frame
whose Content-Security-Policy blocks every script and fetch, and any other is
written into a document without a browsing context, where the full parser runs
and nothing is fetched or run; a page load gives the same tree for the
\dataHtmlPageChecked of them that each run samples. \texttt{DOMParser} would
not: Chromium's and WebKit's read some of them otherwise.
Each browser judges its \dataHtmlBrowserCases cases against its own reading of
the benign value. The browsers give parse5's verdict in all but
\dataHtmlBrowserDiffers cases, all of them concatenation in an \texttt{option},
where parse5 drops markup that the browsers keep.

\paragraph{Escaping}
The per-context arm applies the four rules of the OWASP cheat
sheet~\cite{owasp_xss}: entities in element text, \verb|&#xHH;| for every
non-alphanumeric character below 256 in an attribute value, \verb|\xHH| in the
JavaScript string, and \verb|\HHHHHH| in CSS. \texttt{html/template} does not
take the \texttt{onclick} host, where it
writes the value as a JavaScript string literal, not as the attribute's value.
Neither admits a breakout in the hosts that are not semantic sinks
(\dataHtmlEscapeContextBreakout, \dataHtmlAutoescapeBreakout). Escaped for
JavaScript or CSS, the value in the hole differs from $v$, so those cases count
as altered. Sanitizers such as
DOMPurify keep the markup they allow, for rich HTML outside the claim
(\cref{sec:resist:xss}), and Trusted Types admit a value by its type
(\cref{sec:related}); neither delimits a value, so neither is an arm.

\paragraph{JavaScript string}
The script host writes \verb|var x = "|$v$\verb|";|. Its JavaScript verdict
reads the string literal by the ECMAScript grammar: inert if the literal ends at
the quote that the host wrote and holds $v$; otherwise the script is compiled,
never run, and counts as a breakout if it compiles and as malformed if not. Concatenation breaks out of the string in
\dataHtmlJsConcatBreakout cases. The per-context arm and \texttt{html/template}
hold \dataHtmlJsEscapeContextInert and \dataHtmlJsAutoescapeInert strings, the
JavaScript-encoded hole all \dataHtmlJsComposeInert of \dataHtmlJsComposeCases.

\paragraph{Mutation XSS}
A hole holds its value as text, so with matchertext the only mutation-XSS route
for a value is the second parse after a serializer writes that text into raw text
(L2). DOMPurify's mXSS fixtures~\cite{dompurify_fixtures}
(\dataHtmlMxssProbes values) take that route through every arm as values: the
stock writer lets \dataHtmlMxssStockBreakout cases out, and no hole, guarded
writer, or matchertext writer lets any out, in parse5 or a browser
(\dataHtmlMxssContainedBreakout of \dataHtmlMxssCases cases). Their other route,
markup that moves to another namespace when a parsed fragment is written and
parsed again, needs the value parsed as markup, which no hole does; that is the
rich HTML of \cref{sec:resist:xss}, outside the claim. The default guard checks
raw text and comments and does not cover it either.

\paragraph{Client path and fail-stop}
DOM-based XSS (\dataSubTechniqueDomBased CVEs, named so or with a DOM source)
runs in the page, so the page composes the HTML: the client arm runs
\texttt{composeHtml}, serializes the \texttt{html} element with the guard, and
sets it by \texttt{innerHTML}, which parses it as a fragment. The three browsers
read no breakout from its \dataHtmlComposeDomCases documents. Browsers parse
incrementally and recover from every error, so the matchertext checks fail stop
at the writer: \textsc{Verify}, \texttt{composeHtml}, and the guard refuse before
any HTML leaves, and a stock browser receives HTML without holes. A reader with
the option stops at a malformed hole, as an XML parser stops at a well-formedness
error~\cite{xml10}, which browsers do for XHTML. An element or attribute hole
lies inside one tag token, which a streaming tokenizer holds until its \verb|>|
anyway, and an MDATA section ends at matcher balance, with no look-ahead.
